% Generated draft: requires mathematical and source review.
\documentclass[11pt]{article}
\usepackage[margin=25mm]{geometry}
\usepackage{fontspec}
\setmainfont{Latin Modern Roman}
\usepackage{amsmath,amssymb,mathtools,unicode-math}
\setmathfont{Latin Modern Math}
\usepackage{enumitem,needspace,xurl,hyperref}
\setlength{\parindent}{0pt}
\setlength{\parskip}{5pt}
\setlength{\emergencystretch}{3em}
\newcommand{\N}{\mathbb{N}}
\newcommand{\Z}{\mathbb{Z}}
\newcommand{\Q}{\mathbb{Q}}
\newcommand{\R}{\mathbb{R}}
\newcommand{\C}{\mathbb{C}}
\newcommand{\F}{\mathbb{F}}
\newcommand{\A}{\mathbb{A}}
\newcommand{\PP}{\mathbb{P}}
\newcommand{\E}{\mathbb{E}}
\newcommand{\eps}{\varepsilon}
\newcommand{\dd}{\,\mathrm{d}}
\begin{document}
\section*{UnsolvedMath ID: 1244}
\textbf{Source problem code:} catalogue entry
\subsection{Definitions and mathematical statement}
For a real number $y$, let $\lfloor y\rfloor$ be its greatest integer not exceeding $y$, and let $\{y\}=y-\lfloor y\rfloor\in[0,1)$ be its fractional part. The literal question in the supplied statement is
\[\exists x\in\R\ \ \forall n\in\Z_{\geq1},\qquad \left\{x\left(\frac32\right)^n\right\}<\frac12.\]
The strict inequality and the starting index $n=1$ are retained. No condition $x>0$ is present in this literal statement.

The supplied literature triage checked 2026-08-17 records the affirmative witness $x=0$, with all displayed fractional parts equal to zero, and therefore labels this exact record solved / SOLVED-IN-LITERATURE. That source-provided observation must not be reported as a resolution of Mahler's standard positive-real problem. The corrected existence question identified by the review adds $x>0$ to the same fractional-part condition; the background attributes to Kurt Mahler in the 1960s the conjectural nonexistence of such positive real numbers, called $Z$-numbers in the supplied references.

For the positive-real problem the review reports Mahler's result that at most one $Z$-number lies in each unit interval and later scarcity and distribution restrictions; it reports the standard existence question as still open. These are supplied dated assessments. The registry's separate status field is open, reflecting a conflict between the broad historical label and the unrestricted literal wording; the two targets are kept distinct.
\subsection{Short English statement}
The stored question allows every real $x$ and is reported answered affirmatively by $x=0$. Mahler's intended problem requires $x>0$; the supplied review says that positive-real existence question remains open. The omitted positivity condition is therefore essential.
\subsection{Sources}
\begin{itemize}
\item Local UnsolvedMath registry, permanent ID 1244, source code catalogue entry, ``Mahler's 3/2 Problem''; supplied statement, historical background, and corrective literature triage.
\item K. Mahler, ``An unsolved problem on the powers of 3/2,'' \emph{Journal of the Australian Mathematical Society} 8 (1968), 313--321: \url{https://doi.org/10.1017/S144678870000613X}. Supplied primary source for the positive-real problem.
\item A. Dubickas, ``On the fractional parts of rational powers,'' \emph{Acta Arithmetica} 141 (2010), 103--110: \url{https://doi.org/10.4064/aa141-2-1}. The supplied review says it defines a $Z$-number with positive parameter and treats nonexistence as Mahler's conjecture.
\item S. Akiyama et al., ``Rational self-affine tiles,'' \emph{American Mathematical Monthly} 129 (2022), 537--555: \url{https://doi.org/10.1080/00029890.2022.2061281}. Supplied specialist secondary exposition reporting the positive-real question open.
\item Supplied literature triage checked 2026-08-17, SOLVED-IN-LITERATURE for the unrestricted exact wording only; it explicitly prohibits treating this as a solution to the positive-real problem and cautions that dated triage does not exclude later or unindexed work.
\end{itemize}
\end{document}
